\documentclass[12pt,a4paper]{article}
\usepackage[a4paper,margin=0.85in]{geometry}
\usepackage[T1]{fontenc}
\usepackage{mathptmx}
\usepackage{graphicx}
\usepackage{array,tabularx,longtable,multirow}
\usepackage{pdflscape}
\usepackage{setspace}
\usepackage{ragged2e}
\usepackage{fancyhdr}
\usepackage{enumitem}
\usepackage[hidelinks]{hyperref}
\setstretch{1.15}
\setlength{\parindent}{0pt}
\setlength{\parskip}{0.45em}
\pagestyle{fancy}
\fancyhf{}
\fancyfoot[C]{\thepage}
\renewcommand{\headrulewidth}{0pt}
\newcommand{\ProjectTitle}{ProjectGuard: Academic Project Management and Similarity Screening}
\newcommand{\SupervisorName}{\rule{6cm}{0.4pt}}
\newcommand{\MemberOne}{\rule{6cm}{0.4pt}}
\newcommand{\MemberTwo}{\rule{6cm}{0.4pt}}
\newcommand{\MemberThree}{\rule{6cm}{0.4pt}}
\newcommand{\MemberFour}{\rule{6cm}{0.4pt}}
\newcommand{\fillline}{\rule{0.95\linewidth}{0.4pt}}
\newcommand{\progressrow}[4]{#1 & & & & & & #2 & #3 & #4 & \\[0.8em]\hline}
\begin{document}

% Replace the macros above with actual project and participant details.
\begin{titlepage}
\centering
\vspace*{0.7cm}
{\bfseries\LARGE MINI-PROJECT LOGBOOK\par}
\vspace{1.2cm}
{\bfseries\large GROUP MEMBERS\par}
\vspace{0.5cm}
\begin{tabular}{@{}cl@{}}
1. & \MemberOne\\[0.5cm]
2. & \MemberTwo\\[0.5cm]
3. & \MemberThree\\[0.5cm]
4. & \MemberFour
\end{tabular}
\vspace{1.2cm}

{\bfseries Project Title\par}
\ProjectTitle\par
\textbf{Project status: Completed and deployed}\par
\vspace{0.9cm}
{\bfseries Supervisor\par}
Prof. \SupervisorName\par
\vfill
\includegraphics[width=2.8cm]{institute_logo.jpeg}\par
\vspace{0.5cm}
{\bfseries Department of Artificial Intelligence and Machine Learning\par}
SIES Graduate School of Technology (Autonomous)\par
Affiliated to University of Mumbai\par
Nerul, Navi Mumbai -- 400 706\par
\vspace{0.4cm}
(AY 2026--27)
\end{titlepage}

\section*{INSTITUTE VISION \& MISSION}
\subsection*{VISION}
\begin{quote}
To be a centre of excellence in Education and Technology committed towards socio-economic advancement of the country.
\end{quote}
\subsection*{MISSION}
\begin{enumerate}[leftmargin=1.6em,itemsep=0.5em]
\item To impart advanced knowledge in Engineering and Technology.
\item To transform young minds towards professional competence by inculcating values and developing skills.
\item To promote research and ensure continuous value addition among students and employees.
\item To strengthen association with industry, research organizations and alumni to enhance knowledge on current technologies.
\item To promote next generation technocracy and nurture entrepreneurial culture for socio-economic growth.
\end{enumerate}
\clearpage

\section*{ARTIFICIAL INTELLIGENCE AND MACHINE LEARNING DEPARTMENT}
\subsection*{DEPARTMENT VISION}
Impart prime content education to develop technically proficient and socially responsible engineers in the emerging branches of AI.
\subsection*{DEPARTMENT MISSION}
\begin{description}[leftmargin=0.9cm,style=nextline,itemsep=0.25em]
\item[M1] To provide students with a comprehensive education to prepare them for the rigors in the domain of AI.
\item[M2] To establish a platform for academics, innovation, professional growth, and social interaction.
\item[M3] To encourage the capacity for lifelong learning necessary for career advancement.
\item[M4] To help build relationships with alumni and industry for the betterment of students' overall growth.
\end{description}
\subsection*{PROGRAMME EDUCATIONAL OBJECTIVES (PEOs)}
\begin{description}[leftmargin=1.1cm,style=nextline,itemsep=0.2em]
\item[PEO1] Develop leadership abilities for professional advancement.
\item[PEO2] Pursue higher education to learn and advance their careers.
\item[PEO3] Promote services in the field of Artificial Intelligence Engineering through excellent technical and communication skills.
\end{description}
\clearpage

\section*{PROGRAMME OUTCOMES (POs)}
\small
\begin{longtable}{|p{0.10\textwidth}|p{0.78\textwidth}|}
\hline \textbf{PO} & \textbf{Outcome}\\\hline
\endfirsthead
\hline \textbf{PO} & \textbf{Outcome}\\\hline
\endhead
PO1 & \textbf{Engineering Knowledge:} Apply knowledge of mathematics, natural science, computing, engineering fundamentals and an engineering specialization to develop the solution of complex engineering problems.\\\hline
PO2 & \textbf{Problem Analysis:} Identify, formulate, review research literature and analyze complex engineering problems reaching substantiated conclusions with consideration for sustainable development.\\\hline
PO3 & \textbf{Design/Development of Solutions:} Design creative solutions for complex engineering problems and design or develop systems, components, and processes to meet identified needs with consideration for public health and safety, whole-life cost, net zero carbon, culture, society and environment as required.\\\hline
PO4 & \textbf{Conduct Investigations of Complex Problems:} Conduct investigations using research-based knowledge, including design of experiments, modelling, analysis and interpretation of data to provide valid conclusions.\\\hline
PO5 & \textbf{Engineering Tool Usage:} Create, select and apply appropriate techniques, resources and modern engineering and IT tools, including prediction and modelling, recognizing their limitations.\\\hline
PO6 & \textbf{The Engineer and The World:} Analyze and evaluate societal and environmental aspects while solving complex engineering problems for their impact on sustainability with reference to economy, health, safety, legal framework, culture and environment.\\\hline
PO7 & \textbf{Ethics:} Apply ethical principles and commit to professional ethics, human values, diversity and inclusion; adhere to national and international laws.\\\hline
PO8 & \textbf{Individual and Collaborative Team Work:} Function effectively as an individual and as a member or leader in diverse or multidisciplinary teams.\\\hline
PO9 & \textbf{Communication:} Communicate effectively and inclusively within the engineering community and society at large, including effective reports, design documentation and presentations that consider cultural, language and learning differences.\\\hline
PO10 & \textbf{Project Management and Finance:} Apply engineering management principles and economic decision-making to individual work, team leadership and multidisciplinary projects.\\\hline
PO11 & \textbf{Life-Long Learning:} Recognize the need and preparation for independent lifelong learning, adaptability to emerging technologies and critical thinking in technological change.\\\hline
\end{longtable}
\normalsize

\section*{PROGRAMME SPECIFIC OUTCOMES (PSOs)}
\textbf{PSO1:} Students are exposed to cutting-edge trends and technology to prepare them for challenging work in industry, higher education and professional careers.

\textbf{PSO2:} Practice skills learnt in Artificial Intelligence, Data Mining, Machine Learning, Cloud Computing, Optimization, Networking and Security, and Deep Learning in domains such as Healthcare, Education, Agriculture, Finance and Social Media.
\clearpage

\section*{STUDENT INFORMATION}
\textbf{Project Title:} \ProjectTitle

\vspace{0.5cm}
\small
\renewcommand{\arraystretch}{2.2}
\begin{tabularx}{\textwidth}{|p{0.20\textwidth}|*{4}{>{\centering\arraybackslash}X|}}
\hline
 & \textbf{Student 1} & \textbf{Student 2} & \textbf{Student 3} & \textbf{Student 4}\\\hline
Roll No. & & & &\\\hline
Name & & & &\\\hline
Class with Division & & & &\\\hline
Contact No. & & & &\\\hline
E-mail & & & &\\\hline
\end{tabularx}
\normalsize
\renewcommand{\arraystretch}{1}
\vspace{1cm}
\section*{INSTRUCTIONS TO STUDENTS}
\begin{enumerate}[leftmargin=1.6em]
\item The logbook must be submitted to the Supervisor for verification and evaluation of project activities at least once a week.
\item The logbook, duly signed by the Supervisor, must be submitted with the project report for evaluation at the end of the semester to the department.
\end{enumerate}
\clearpage

\section*{LETTER OF ACCEPTANCE}
\vspace{0.7cm}
I, the undersigned, Prof. \SupervisorName, working in the Department of Artificial Intelligence and Machine Learning, am willing to guide the project titled:

\vspace{0.4cm}
\begin{center}\textbf{\ProjectTitle}\end{center}
\vspace{0.4cm}
for Mini Project 2, Semester V, for the academic year 2026--27. The names of the students are:
\begin{enumerate}[leftmargin=1.8em,itemsep=0.5em]
\item \MemberOne
\item \MemberTwo
\item \MemberThree
\item \MemberFour
\end{enumerate}
\vfill
\begin{tabularx}{\textwidth}{>{\centering\arraybackslash}X>{\centering\arraybackslash}X>{\centering\arraybackslash}X}
\rule{0.8\linewidth}{0.4pt} & \rule{0.8\linewidth}{0.4pt} & \rule{0.8\linewidth}{0.4pt}\\
(Name) Mini-Project Guide & Prof. Jegan Nadar\\Mini Project Coordinator & Dr. Kalyani Pampattiwar\\I/C Head of Department
\end{tabularx}
\clearpage

\section*{COURSE OUTCOMES}
\footnotesize
\begin{tabularx}{\textwidth}{|p{0.09\textwidth}|X|p{0.16\textwidth}|p{0.12\textwidth}|}
\hline
\textbf{CO} & \textbf{Course Outcome} & \textbf{POs covered} & \textbf{PSOs covered}\\\hline
CO1 & Identify and define real-world or research problems through requirement gathering and stakeholder interaction. & PO2, PO4, PO6, PO9, PO11 & PSO1\\\hline
CO2 & Develop project objectives, scope and feasibility using software engineering and SDLC principles. & PO1, PO2, PO3, PO5, PO10, PO11 & PSO1\\\hline
CO3 & Analyze requirements and develop system designs using standard modelling techniques. & PO1, PO2, PO3, PO4, PO5 & PSO1, PSO2\\\hline
CO4 & Implement solutions using appropriate technologies, engineering practices and version control. & PO1, PO3, PO5, PO11 & PSO1, PSO2\\\hline
CO5 & Test and validate the system through structured test cases, datasets and simulations. & PO2, PO4, PO5 & PSO1, PSO2\\\hline
CO6 & Compose technical documentation, present project outcomes and demonstrate teamwork and ethics. & PO7, PO8, PO9 & PSO1\\\hline
\end{tabularx}
\normalsize
\clearpage

\begin{landscape}
\section*{CO--PO--PSO MAPPING}
\begin{center}
\small
\renewcommand{\arraystretch}{1.9}
\setlength{\tabcolsep}{7pt}
\begin{tabular}{|l|*{11}{c|}*{2}{c|}}
\hline
 & PO1 & PO2 & PO3 & PO4 & PO5 & PO6 & PO7 & PO8 & PO9 & PO10 & PO11 & PSO1 & PSO2\\\hline
CO1 & & 3 & & 3 & & 3 & & & 3 & & & 3 &\\\hline
CO2 & 3 & 3 & 3 & & 3 & & & & & 3 & 3 & 2 &\\\hline
CO3 & 3 & 3 & 3 & 3 & 3 & & & & & & & 3 & 3\\\hline
CO4 & 3 & & 2 & & 3 & & & 2 & & 2 & 3 & 3 & 3\\\hline
CO5 & & 3 & & 3 & 3 & & & & & & & 3 & 3\\\hline
CO6 & & & & & & & 3 & 3 & 3 & & & 3 &\\\hline
\end{tabular}
\end{center}
\end{landscape}

\section*{SCHEDULE / PROGRESS REPORT}
\begin{tabularx}{\textwidth}{|p{0.24\textwidth}|X|}
\hline
Title of the Project & \ProjectTitle\\\hline
Group No. & \rule{4cm}{0.4pt}\\\hline
Student 1 & \MemberOne\\\hline
Student 2 & \MemberTwo\\\hline
Student 3 & \MemberThree\\\hline
Student 4 & \MemberFour\\\hline
Name of the Supervisor & Prof. \SupervisorName\\\hline
\end{tabularx}
\clearpage

\begin{landscape}
\section*{WEEKLY PROGRESS RECORD}
\small
\renewcommand{\arraystretch}{1.45}
\setlength{\tabcolsep}{3pt}
\begin{longtable}{|c|p{1.6cm}|c|c|c|c|p{6.6cm}|p{3.4cm}|p{1.3cm}|p{2.9cm}|}
\hline
\multirow{2}{*}{Sr. No.} & \multirow{2}{*}{Date} & \multicolumn{4}{c|}{Attendance} & \multirow{2}{*}{List of scheduled tasks} & \multirow{2}{*}{Remark} & Mapping & \multirow{2}{*}{Supervisor signature}\\\cline{3-6}\cline{9-9}
 & & 1 & 2 & 3 & 4 & & & CO &\\\hline
\endfirsthead
\hline
Sr. No. & Date & 1 & 2 & 3 & 4 & List of scheduled tasks & Remark & CO & Supervisor signature\\\hline
\endhead
\progressrow{1}{Defined the project problem, requirements, stakeholders, and scope.}{Completed}{CO1}
\progressrow{2}{Reviewed related project portals and lexical/semantic similarity methods.}{Completed}{CO1}
\progressrow{3}{Designed architecture, data model, API flows, and team ownership rules.}{Completed}{CO2, CO3}
\progressrow{4}{Built student and mentor portal layouts and workflow screens.}{Completed}{CO3, CO4}
\progressrow{5}{Moved active data to PostgreSQL and added pgvector project storage.}{Completed}{CO4}
\progressrow{6}{Integrated Clerk sessions, enrollment codes, and role checks.}{Completed}{CO4, CO5}
\progressrow{7}{Implemented team creation, joining, invitations, and submissions.}{Completed}{CO4}
\progressrow{8}{Implemented mentor roster, assignment, reviews, and admin operations.}{Completed}{CO4}
\progressrow{9}{Added PDF, DOCX, and TXT validation, extraction, and limits.}{Completed}{CO4, CO5}
\progressrow{10}{Integrated TF--IDF, Qwen embeddings, hybrid scores, and vector search.}{Completed}{CO4, CO5}
\progressrow{11}{Tested authorization, team ownership, upload errors, and score responses.}{Completed}{CO5}
\progressrow{12}{Configured frontend/backend build, HTTPS origins, secrets, and database.}{Completed}{CO4, CO5}
\progressrow{13}{Completed integration checks, documentation, deployment, and handover.}{Completed}{CO5, CO6}
\end{longtable}
\normalsize
\vspace{0.4cm}
\hfill \rule{6cm}{0.4pt}\\
\hfill Signature of the Supervisor
\end{landscape}

\section*{PROJECT COMPLETION SUMMARY}
ProjectGuard was completed as an integrated academic project management and similarity-screening system. The delivered work includes student and mentor workflows, administrator enrollment, PostgreSQL and pgvector persistence, document extraction, TF--IDF and Qwen-based comparison, authentication checks, and deployment configuration. The similarity score is used for human-guided review; a numerical detection-accuracy result is not entered in this logbook without a labeled evaluation record.

The progress entries above describe completed work. Enter the actual meeting dates, attendance, supervisor remarks, and signatures from the team's records before submitting the signed logbook.
\clearpage

\section*{EXAMINER'S FEEDBACK FORM}
\begingroup
\setstretch{1}
\begin{tabularx}{\textwidth}{p{0.38\textwidth}X}
Name of External Examiner: & \fillline\\[0.5em]
College of External Examiner: & \fillline\\[0.5em]
Name of Internal Examiner: & \fillline\\[0.5em]
Date of Examination: & \fillline\\[0.5em]
No. of Students in Project Team: & \fillline\\[0.5em]
Separate Lab Available: & Yes \quad / \quad No
\end{tabularx}
\par\vspace{0.35cm}

\textbf{Student Performance Analysis}\\
\small Tick one rating for each observation.

\vspace{0.2cm}
\small
\renewcommand{\arraystretch}{1.35}
\setlength{\tabcolsep}{4pt}
\begin{tabular}{|>{\centering\arraybackslash}p{0.55cm}|p{8.1cm}|>{\centering\arraybackslash}p{1.45cm}|>{\centering\arraybackslash}p{1.45cm}|>{\centering\arraybackslash}p{1.45cm}|}
\hline
\textbf{No.} & \textbf{Observation} & \shortstack{\textbf{Excellent}\\\textbf{(3)}} & \shortstack{\textbf{Very Good}\\\textbf{(2)}} & \shortstack{\textbf{Good}\\\textbf{(1)}}\\\hline
1 & Quality of problem and clarity & & &\\\hline
2 & Innovativeness in solutions & & &\\\hline
3 & Cost effectiveness and societal impact & & &\\\hline
4 & Full functioning of working model as per stated requirements & & &\\\hline
5 & Effective use of skill sets & & &\\\hline
6 & Effective use of standard engineering norms & & &\\\hline
7 & Contribution of an individual as member or leader & & &\\\hline
8 & Clarity in written and oral communication & & &\\\hline
9 & Overall performance & & &\\\hline
\end{tabular}
\par
\normalsize
\renewcommand{\arraystretch}{1}
\vspace{0.35cm}

Can the same mini project be extended to the next semester by adding new objectives or ideas? \quad Yes \quad / \quad No

If yes, suggest a new innovative technique, idea or objective related to this project:

\vspace{0.3cm}\fillline

\vspace{0.45cm}\fillline

\vfill
\begin{tabularx}{\textwidth}{>{\centering\arraybackslash}X>{\centering\arraybackslash}X}
\rule{0.7\linewidth}{0.4pt} & \rule{0.7\linewidth}{0.4pt}\\
Signature of External Examiner & Signature of Internal Examiner
\end{tabularx}
\endgroup
\end{document}
